\chapter{Conclusion}
\thispagestyle{empty}

\section{Final System}

Based on the results presented in this thesis, the Thorlabs diode laser together with the camera-based detection system is the most suitable choice for integration into the final setup.

The camera-based detection method was less affected by environmental noise. The homodyne quadrature interferometer successfully demonstrated the feasibility of direction-sensitive displacement measurements. However, since the quadrature setup relies on two photodiodes, whose signals were affected by fluctuations and noise, the camera-based detection method is expected to provide more reliable results for future measurements.


\section{Outlook}

The interferometric displacement measurement system developed in this thesis provides a solid basis for future work. Although the system was successfully implemented, several aspects could be further improved.

One important step would be to increase the stability of the experimental setup. During the measurements, environmental influences such as vibrations affected the interference signal. A better isolated setup would reduce these disturbances and improve the measurement reliability.

In addition, the origin of the day-to-day variations observed for the diode laser should be investigated. Identifying and eliminating the cause of these fluctuations could further improve the stability of the system.

Finally, the complete measurement system should be tested in the intended experimental setup. 


